\begin{abstract}
The ST-CC framework discovers mobility patterns by jointly clustering the regions and time slots of a count matrix into a regular grid, but every split applies to the whole matrix, so interaction that occurs only inside a block remains hidden and there is no statistical criterion to decide whether a block should be subdivided. This thesis proposes AST-CC, a tree that refines the ST-CC partition block by block through hypothesis tests derived from the I-divergence, calibrated with a day bootstrap that preserves the real variability between days, adjusted for the search of the best split, and with the significance level allocated along the tree. An effect-size threshold separates relevant structure from structure that is merely detectable because of the large volume of data, and among the admissible splits the one that achieves its gain with fewer parameters is chosen. On synthetic data, the method keeps the rate of spurious refinements below 5\,\%, whereas a Poisson null produces them in every replicate, almost fully recovers the true structure when the base partition groups poorly, and outperforms finer global grids in every replicate while using fewer parameters. On the Manhattan taxi trips of November 2017, the tree removes 95.9 and 90.7\,\% of the interaction that ST-CC leaves inside its blocks in the weekday and holiday periods, outperforms optimized global grids in the weekday period, and in the holiday period loses to them with the ST-CC partition but ties when it starts from a better one. Its splits replicate in the rest of 2017 and reveal timing differences between zones of similar volume, such as the commuting lag and the contrast between nightlife and residential zones.
\end{abstract}
